\documentclass[uplatex,dvipdfmx,a4paper,11pt]{jsarticle}
\input{../../../templates/note.tex}
\noteprint{量子力学1　第4章「数学的準備：デルタ関数・フーリエ級数・フーリエ変換」}{2026-11-25}
%% note-moto	problems\量子力学1\shou04\量子力学1-4.tex
\begin{document}
\begin{kirokuichiran}
  \ichirangyou{\hyperlink{bt-m-019-fourier}{1}}{L1・計算}{◎◎}{S}{2026-10-18}
  \ichirangyou{\hyperlink{bt-m-028-delta}{2}}{L1・計算}{◎○◎}{A}{2026-10-28 (2)記述}
  \ichirangyou{3}{L1・計算}{}{未着手}{}
  \ichirangyou{\hyperlink{bt-m-032-deltaprop}{4}}{L2・導出}{○×××}{E}{2026-10-28 (1)記述 (2)未着手 (3)未着手 (4)未着手}
  \ichirangyou{5}{L2・計算}{}{未着手}{}
  \ichirangyou{6}{L2・計算}{}{未着手}{}
  \ichirangyou{7}{L2・計算}{}{未着手}{}
  \ichirangyou{8}{L3・導出}{}{未着手}{}
  \ichirangyou{9}{L3・計算}{}{未着手}{}
  \ichirangyou{10}{L3・計算}{}{未着手}{}
  \ichirangyou{11}{L4・導出}{}{未着手}{}
\end{kirokuichiran}

\mondaihen
\level{1}
% toubun: kiroku\toubun\量子力学1\019_2026-10-18.tex
\mondaibangou{1}
\begin{mondai}{019-fourier}{計算}{フーリエ変換の定義（係数 $1/\sqrt{2\pi}$ の規約）と，関数の幅と波数の広がり}
フーリエ変換は要点整理の規約で定める。$a>0$ として，関数
  \[
    f(x)=
    \begin{cases}
      1/\sqrt{2a} & (\abs{x}\le a)\\
      0 & (\abs{x}>a)
    \end{cases}
  \]
  を考える。
  \begin{shoumon}
    \item $f(x)$ のフーリエ変換 $F(k)$ を，$k\neq0$ と $k=0$ の場合に分けて求めよ。
    \item $F(k)=0$ となる正の $k$ のうち最小のもの $k_{1}$ を求めよ。また，$f(x)$ が0でない区間の幅 $2a$ と $k_{1}$ の積が $a$ によらないことを示し，$a$ を大きくすると $F(k)$ の形がどう変わるかを述べよ。
  \end{shoumon}
\notekiroku{2026-10-18 S ◎◎}
\begin{kotaeutsushi}{2026-10-18}
\toi{1}{◎}{}
i) $k\neq0$ のとき
\begin{align*}
F(k)&=\frac{1}{\sqrt{2\pi}}\int\dd x\,\ee^{-\ii kx}f(x)\\
&=\frac{1}{\sqrt{2\pi}}\int_{-a}^{a}\kakimachigai{\frac{\ee^{-\ii kx}}{\sqrt{2a}}}{\frac{\ee^{-\ii kx}}{\sqrt{2a}}\dd x}\\
&=\frac{1}{2\sqrt{\pi a}}\left[\frac{\ee^{-\ii kx}}{-\ii k}\right]_{-a}^{a}\\
&=\frac{\ii}{2k\sqrt{\pi a}}\left(\ee^{-\ii ka}-\ee^{\ii ka}\right)\\
&=\frac{\sin(ka)}{k\sqrt{\pi a}}
\end{align*}
ii) $k=0$ のとき
\begin{align*}
F(0)&=\frac{1}{\sqrt{2\pi}}\int\dd x\,f(x)\\
&=\frac{1}{2\sqrt{\pi a}}\int_{-a}^{a}\dd x\\
&=\sqrt{\frac{a}{\pi}}
\end{align*}
\toi{2}{◎}{}
\[ \frac{\sin(ka)}{k\sqrt{\pi a}}=0\qquad\therefore\ \sin(ka)=0 \]
\begin{align*}
&\therefore\ ka=n\pi\quad(n=1,2,3)\\
&\therefore\ k_{1}=\frac{\pi}{a}
\end{align*}
\[ 2ak_{1}=2\pi \]
$a$ を大きくすると，$k_{1}$ は小さくなるため，関数が鋭くなる。

（図：$k=0$ の山の両側に小さな振動が続く $F(k)$ の形が，矢印の先では幅の狭い山になる）
\end{kotaeutsushi}
\end{mondai}

% toubun: kiroku\toubun\量子力学1\028_2026-10-28.tex
\mondaibangou{2}
\begin{mondai}{028-delta}{計算}{デルタ関数の定義の使い方，階段関数の微分}
次の積分を求めよ。$\theta(x)$ は $x\ge 0$ で 1，$x<0$ で 0 の階段関数とする。
  \begin{shoumon}
    \item 次の積分。
    \[ \int_{-\infty}^{\infty}(x^{2}+3x)\,\delta(x-2)\dd x \]
    \item 次の積分。
    \[ \int_{0}^{\infty}\ee^{-x}\,\delta(x+1)\dd x \]
    \item 次の積分を，$\delta(x)=\dd \theta/\dd x$ を使わずに部分積分で求めよ（表面項の値も確かめること）。また，その結果が $\delta(x-1)$ を使って計算した値と一致することを確かめよ。
    \[ \int_{-\infty}^{\infty}\ee^{-x^{2}}\,\dv{\theta(x-1)}{x}\dd x \]
  \end{shoumon}
\notekiroku{2026-10-28 A ◎○◎}
\begin{kotaeutsushi}{2026-10-28}
\toi{1}{◎}{}
\[ 2^{2}+3\cdot2=10 \]
\toi{2}{○}{記述}
\[ 0 \]
\nuke{$x=-1$ は区間外}
\naze{答えの 0 は正しいが根拠が無い。$\delta(x+1)$ が値をもつ $x=-1$ は積分区間 $0\le x<\infty$ の外にあるので積分は 0 になる，の一言が要る。}
\toi{3}{◎}{}
\begin{align*}
&\int_{-\infty}^{\infty}\ee^{-x^{2}}\dv{\theta(x-1)}{x}\dd x\\
&=\underbrace{\Bigl[\ee^{-x^{2}}\theta(x-1)\Bigr]_{-\infty}^{\infty}}_{0}\\
&\qquad+\int_{-\infty}^{\infty}2x\ee^{-x^{2}}\kakimachigai{\theta(x)}{\theta(x-1)}\dd x\\
&=\int_{1}^{\infty}2x\ee^{-x^{2}}\dd x
\end{align*}
（脇で置換：$\ee^{-x^{2}}=z$，$-2x\ee^{-x^{2}}\dd x=\dd z$，$x:1\to\infty$ で $z:1/\ee\to0$）
\[ =\int_{0}^{1/\ee}\dd z=\frac{1}{\ee} \]
\[ \int_{-\infty}^{\infty}\ee^{-x^{2}}\kakimachigai{\delta(x-1)}{\delta(x-1)\dd x}=\frac{1}{\ee} \]
\end{kotaeutsushi}
\begin{kizuki}
\item デルタ関数の積分では，まず $\delta$ の立つ点（$\delta(x+1)$ なら $x=-1$）が積分区間に入っているかを書く。入っていなければ 0，入っていればその点での値になる。答えが 0 のときは，この一言がそのまま答えの根拠になる。
\end{kizuki}
\begin{tadashii}
\houshin デルタ関数の定義 $\int f(x)\delta(x-a)\dd x=f(a)$ を，$x=a$ が積分区間に入っているかに注意して使う。(3) は部分積分で微分を $\ee^{-x^{2}}$ に移す。
\keisan \shoumonkai{1} $x=2$ は積分区間に入っているので，$f(x)=x^{2}+3x$ の $x=2$ での値になる。
\[ \int_{-\infty}^{\infty}(x^{2}+3x)\delta(x-2)\dd x=2^{2}+3\cdot 2=10 \]

\shoumonkai{2} $\delta(x+1)=\delta(x-(-1))$ は $x=-1$ にだけ寄与をもつ。積分区間 $0\le x<\infty$ は $x=-1$ を含まないので，積分は 0 である。

\shoumonkai{3} 部分積分で微分を $\ee^{-x^{2}}$ に移す。
\begin{align*}
\int_{-\infty}^{\infty}\ee^{-x^{2}}\dv{\theta(x-1)}{x}\dd x&=\Bigl[\ee^{-x^{2}}\theta(x-1)\Bigr]_{-\infty}^{\infty}-\int_{-\infty}^{\infty}\bigl(-2x\ee^{-x^{2}}\bigr)\theta(x-1)\dd x
\end{align*}
表面項は，$x\to\infty$ で $\ee^{-x^{2}}\to 0$，$x\to-\infty$ で $\theta(x-1)=0$ なので 0。第2項は $\theta(x-1)$ が $x\ge 1$ でだけ 1 なので，
\[ \int_{1}^{\infty}2x\ee^{-x^{2}}\dd x=\Bigl[-\ee^{-x^{2}}\Bigr]_{1}^{\infty}=\ee^{-1} \]
一方，$\dd \theta(x-1)/\dd x=\delta(x-1)$ を使えば，定義から $\ee^{-1^{2}}=\ee^{-1}$ となり，一致する。
\kotae{\shoumonkai{1} $10$ \shoumonkai{2} $0$（$x=-1$ が積分区間の外） \shoumonkai{3} $\ee^{-1}$（$\delta(x-1)$ を使った値と一致）}
\begin{ronten}
  \item \shoumonkai{2} デルタ関数の立つ点 $x=-1$ が積分区間 $0\le x<\infty$ の外にあるので，積分が 0 になることを述べる。
  \item \shoumonkai{3} 部分積分の表面項が 0 になる理由を，$x\to\infty$ と $x\to-\infty$ のそれぞれについて書く。
  \item \shoumonkai{3} $\theta(x-1)$ が $x\ge1$ でだけ 1 なので，残る積分の区間が $x\ge1$ に狭まることを書く。
\end{ronten}
\kenzan (3) の値 $\ee^{-1}\approx 0.368$ は $0<\ee^{-1}<1$ で，$\ee^{-x^{2}}$ の最大値 1 を超えない。(2) は積分区間を $-\infty$ からに広げると $\ee^{1}$ になり，区間の端が効いていることが分かる。
\tsumazuki (2) で区間を見ずに $\ee^{-(-1)}=\ee$ と答える誤りが多い。デルタ関数は「立つ点」が積分区間の中にあるときだけ値を拾う。
\end{tadashii}
\end{mondai}

\level{2}
% toubun: kiroku\toubun\量子力学1\032_2026-10-28.tex
\mondaibangou{4}
\begin{mondai}{032-deltaprop}{導出}{変数変換によるデルタ関数の性質，デルタ関数の微分}
デルタ関数の定義 $\int f(x)\delta(x-a)\dd x=f(a)$ と，部分積分が使えて表面項が消えることだけを使って，次を示せ。$f$ は滑らかな関数とする。
  \begin{shoumon}
    \item 0 でない実数 $c$ に対して $\delta(cx)=\delta(x)/\abs{c}$（両辺に $f(x)$ を掛けて積分したものが等しいという意味）。
    \item $a>0$ に対して次の式。$x>0$ と $x<0$ に分けて，それぞれで $y=x^{2}-a^{2}$ と置換せよ。
    \[ \delta(x^{2}-a^{2})=\frac{\delta(x-a)+\delta(x+a)}{2a} \]
    \item 次の式（$\delta'$ は $\delta$ の導関数，$b$ は実数）。
    \[ \int_{-\infty}^{\infty}f(x)\,\delta'(x-b)\dd x=-f'(b) \]
    \item (2)，(3) を使って，次の2つの積分を求めよ。
    \[ I_{1}=\int_{-\infty}^{\infty}(x^{3}+1)\,\delta(x^{2}-4)\dd x,\qquad I_{2}=\int_{-\infty}^{\infty}x^{3}\,\delta'(x-2)\dd x \]
  \end{shoumon}
\notekiroku{2026-10-28 E ○×××}
\begin{kotaeutsushi}{2026-10-28}
\toi{1}{○}{記述}
i) $c>0$（$cx=z$ とおく）
\begin{align*}
&\int_{-\infty}^{\infty}f(x)\delta(cx)\dd x\\
&=\int_{-\infty}^{\infty}\ayamari{f(x)}{f(z/c)}\delta(z)\frac{\dd z}{c}
\end{align*}
ii) $c<0$（$cx=z$ とおく）
\begin{align*}
&\int_{-\infty}^{\infty}f(x)\delta(cx)\dd x\\
&=\int_{\infty}^{-\infty}\ayamari{f(x)}{f(z/c)}\delta(z)\frac{\dd z}{c}\\
&=\int_{-\infty}^{\infty}f(x)\delta(z)\frac{\dd z}{-c}
\end{align*}
よって
\[ \delta(cx)=\frac{\delta(x)}{\abs{c}} \]
\naze{置換したら $f$ の引数も $x=z/c$ に直す。$\delta(z)$ で $z=0$ の値をとると，i) は $f(0)/c$，ii) は $-f(0)/c$ で，まとめて $f(0)/\abs{c}$ になる。これが $\int f(x)\delta(x)\dd x/\abs{c}=f(0)/\abs{c}$ と等しいことを書いて，初めて結論が言える。}
\toi{2}{×}{未着手}
\[ \int \]
\naze{積分記号を書いたところで止まっている。$\int f(x)\delta(x^{2}-a^{2})\dd x$ を $x>0$ と $x<0$ に分け，それぞれ $y=x^{2}-a^{2}$ と置換する。$x>0$ では $x=\sqrt{y+a^{2}}$，$\dd x=\dd y/(2\sqrt{y+a^{2}})$ で $y$ は $-a^{2}$ から $\infty$ まで動き，$f(a)/(2a)$ が出る。$x<0$ でも同様に $f(-a)/(2a)$ が出て，和が $\int f(x)\{\delta(x-a)+\delta(x+a)\}\dd x/(2a)$ と等しい。}
\toi{3}{×}{未着手}
\mikai
\toi{4}{×}{未着手}
\mikai
\end{kotaeutsushi}
\begin{kizuki}
\item 置換 $z=cx$ をしたら，被積分関数の中の $x$ もすべて $z/c$ に書き換える（$f(x)\to f(z/c)$）。$\delta(z)$ で $z=0$ を代入して $f(0)/c$ のような値まで出し，右辺 $\delta(x)/\abs{c}$ を掛けて積分した $f(0)/\abs{c}$ と見比べてから結論する。
\item (2) は (1) と同じ「置換して $\delta(y)$ の積分にし，$y=0$ での値を代入する」流れを，$x>0$ と $x<0$ の2つの範囲で1回ずつ行う。(3) は部分積分で微分を $\delta$ から $f$ に移し，表面項が消えることを使う。(4) は (2)(3) の結果に代入するだけなので，(2)(3) の導出で止まっても結果の式を使って先に解ける。
\end{kizuki}
\begin{tadashii}
\houshin デルタ関数を含む式の等号は「任意の $f$ を掛けて積分したものが等しい」という意味なので，積分の中で置換・部分積分をして，定義 $\int f(x)\delta(x-a)\dd x=f(a)$ の形に持ち込む。
\keisan \shoumonkai{1} $y=cx$ と置換する。$c>0$ では $\dd x=\dd y/c$ で積分の向きは変わらないので，
\[ \int_{-\infty}^{\infty}f(x)\delta(cx)\dd x=\frac{1}{c}\int_{-\infty}^{\infty}f\Bigl(\frac{y}{c}\Bigr)\delta(y)\dd y=\frac{f(0)}{c} \]
$c<0$ では $x\to+\infty$ のとき $y\to-\infty$ なので上端と下端が入れ替わり，
\[ \int_{-\infty}^{\infty}f(x)\delta(cx)\dd x=\frac{1}{c}\int_{\infty}^{-\infty}f\Bigl(\frac{y}{c}\Bigr)\delta(y)\dd y=-\frac{f(0)}{c} \]
どちらも $f(0)/\abs{c}$ で，これは $\delta(x)/\abs{c}$ に $f$ を掛けて積分したものに等しい。

\shoumonkai{2} 積分を $x>0$ と $x<0$ に分ける。$x>0$ では $y=x^{2}-a^{2}$ は単調増加で，$x=\sqrt{y+a^{2}}$，$\dd x=\dd y/(2\sqrt{y+a^{2}})$，$y$ は $-a^{2}$ から $\infty$ まで動き，$y=0$ を含む。
\[ \int_{0}^{\infty}f(x)\delta(x^{2}-a^{2})\dd x=\int_{-a^{2}}^{\infty}\frac{f\bigl(\sqrt{y+a^{2}}\bigr)}{2\sqrt{y+a^{2}}}\delta(y)\dd y=\frac{f(a)}{2a} \]
$x<0$ では $x=-\sqrt{y+a^{2}}$，$\dd x=-\dd y/(2\sqrt{y+a^{2}})$ で，$x$ が $-\infty$ から $0$ に動くとき $y$ は $\infty$ から $-a^{2}$ に動く。向きを直すと負号が消えて，
\[ \int_{-\infty}^{0}f(x)\delta(x^{2}-a^{2})\dd x=\int_{-a^{2}}^{\infty}\frac{f\bigl(-\sqrt{y+a^{2}}\bigr)}{2\sqrt{y+a^{2}}}\delta(y)\dd y=\frac{f(-a)}{2a} \]
和は $[f(a)+f(-a)]/(2a)$ で，これは右辺 $[\delta(x-a)+\delta(x+a)]/(2a)$ に $f$ を掛けて積分したものに等しい。

\shoumonkai{3} 部分積分をして，表面項が消えることを使う。
\[ \int_{-\infty}^{\infty}f(x)\delta'(x-b)\dd x=\Bigl[f(x)\delta(x-b)\Bigr]_{-\infty}^{\infty}-\int_{-\infty}^{\infty}f'(x)\delta(x-b)\dd x=-f'(b) \]

\shoumonkai{4} $I_{1}$ は (2) で $a=2$ とした場合なので，$f(x)=x^{3}+1$ として
\[ I_{1}=\frac{f(2)+f(-2)}{2\cdot 2}=\frac{9+(-7)}{4}=\frac{1}{2} \]
$I_{2}$ は (3) で $f(x)=x^{3}$，$b=2$ とした場合で，$f'(x)=3x^{2}$ なので $I_{2}=-3\cdot 2^{2}=-12$。
\kotae{\shoumonkai{1} $c$ の正負で場合分けし，どちらも $f(0)/\abs{c}$ \shoumonkai{2} 2つの範囲の寄与 $f(a)/(2a)$，$f(-a)/(2a)$ の和 \shoumonkai{3} 部分積分で $-f'(b)$ \shoumonkai{4} $I_{1}=1/2$，$I_{2}=-12$}
\begin{ronten}
  \item デルタ関数を含む等式は，任意の $f$ を掛けて積分したものの等式として示すことを述べる。
  \item \shoumonkai{1} $c<0$ では置換で積分の上端と下端が入れ替わるので，$c$ の正負で場合分けする理由を書く。
  \item \shoumonkai{2} $x>0$ と $x<0$ のそれぞれで $y=x^{2}-a^{2}$ が単調であり，$y=0$ が各範囲の積分区間に入っていることと，$x<0$ では $\dd x$ の負号が積分の向きの入れ替えで消えることを述べる。
  \item \shoumonkai{3} 部分積分の表面項が消えること（問題の前提）を明記してから，残りの積分に定義を使う。
\end{ronten}
\kenzan 次元：$\delta(x)$ は（$x$ の次元）$^{-1}$ なので，$\delta(x^{2}-a^{2})$ は（長さ）$^{-2}$ で，右辺の $\delta(x\mp a)/(2a)$ も（長さ）$^{-1}\cdot$（長さ）$^{-1}$ で合う。(2) は，零点 $x=a$ の近くで $x^{2}-a^{2}\approx 2a(x-a)$ と近似して (1) を使った結果 $\delta(x-a)/(2a)$ とも一致する。(3) で $f=x$，$b=0$ とすると $\int x\delta'(x)\dd x=-1=-\int\delta(x)\dd x$ となり，$x\delta'(x)=-\delta(x)$ と合う。
\tsumazuki (1) で $c<0$ のときに積分の上端と下端の入れ替わりを見落とすと $\delta(cx)=\delta(x)/c$ と負の値になってしまう。デルタ関数（を掛けて積分したもの）は負にならない。
\end{tadashii}
\end{mondai}

\noteowari
\end{document}
